\documentclass[11pt]{article}
\usepackage[margin=1in]{geometry}
\usepackage{amsmath,amssymb}
\usepackage{booktabs}
\usepackage{hyperref}
\usepackage{listings}
\usepackage{xcolor}

\title{Independent Replication --- Osher \& Sethian (1988)\\
       Level-Set Method for Fronts Propagating with\\
       Curvature-Dependent Speed}
\author{Automated replication (rick-stevens-ai/replication-project)}
\date{}

\begin{document}
\maketitle

\begin{abstract}
We independently re-implement, in pure NumPy on a laptop CPU, the
level-set (Hamilton--Jacobi) front-propagation method of Osher and
Sethian, \emph{J.\ Comput.\ Phys.} \textbf{79}, 12--49 (1988). Four
core testable claims (C1: constant-speed circle expansion; C2:
mean-curvature circle collapse; C2b: non-convex star smoothing; C3:
automatic topological merger of two disks) are reproduced to within
the discretization error predicted by the paper's own analysis.
Verdict: \textbf{REPLICATED}.
\end{abstract}

\section{Paper summary}
Osher and Sethian capture a propagating front
$\gamma(t) \subset \mathbb{R}^N$ implicitly as the zero level set of a
scalar function $\phi(x,t)$ evolving under the Hamilton--Jacobi PDE
\begin{equation}
\phi_t + F(K)\,|\nabla\phi| = 0,
\qquad
K = -\frac{\phi_{xx}\phi_y^{2} - 2\phi_{xy}\phi_x\phi_y + \phi_{yy}\phi_x^{2}}
        {(\phi_x^{2}+\phi_y^{2})^{3/2}} .
\end{equation}
For the constant-speed / convection part they use the fully-upwind
Godunov flux (their Eq.~3.11); the curvature part uses \emph{central}
differences throughout, exactly as recommended in \S III.C.

\section{Claims tested}
\begin{center}
\begin{tabular}{@{}llll@{}}
\toprule
ID  & Claim                                                    & Type   & Result \\
\midrule
C1  & Circle expands at $dR/dt = 1$ under $F{=}1$              & quant. & pass \\
C2  & Circle collapses at $R(t)=\sqrt{R_0^{2}-2\varepsilon t}$ & quant. & pass \\
C2b & Non-convex star smooths, perimeter monotone decreasing    & qual./quant. & pass \\
C3  & Two disks merge automatically, no re-meshing             & qual./quant. & pass \\
C4  & Higher-order ENO+RK ordering                             & quant. & not tested (out of scope) \\
C5  & 3-D / N-D generalization                                 & qual.  & not tested (out of scope) \\
\bottomrule
\end{tabular}
\end{center}

\section{Method}
Independent NumPy implementation in \texttt{work/levelset.py} on a uniform
Cartesian grid with periodic BCs; upwind Godunov flux for the convection
term, all-central stencils for the curvature term, and forward-Euler time
stepping with the paper's hyperbolic and parabolic CFL bounds.

\section{Results vs paper}
\subsection*{C1 --- Expanding circle (F=1)}
$R_0=0.25$, $T=0.5$, $N=201$: numerical $R(T)=0.7473$ vs exact $0.7500$;
\textbf{relative error 0.36\%}, max trajectory error $2.87\times10^{-3}$.

\subsection*{C2 --- Mean-curvature circle collapse}
$\varepsilon=0.10$, $R_0=0.30$, grid refinement study:
\begin{center}
\begin{tabular}{@{}rrrrrr@{}}
\toprule
$N$ & $\Delta x$ & $L^{2}$ err & $L^{\infty}$ err & order $L^{2}$ & order $L^{\infty}$ \\
\midrule
101 & 0.0120 & $9.91{\times}10^{-4}$ & $3.88{\times}10^{-3}$ & --   & --   \\
201 & 0.0060 & $3.07{\times}10^{-4}$ & $1.54{\times}10^{-3}$ & 1.69 & 1.34 \\
301 & 0.0040 & $1.59{\times}10^{-4}$ & $7.01{\times}10^{-4}$ & 1.63 & 1.93 \\
\bottomrule
\end{tabular}
\end{center}
Observed order 1.3--1.9 is consistent with the paper's discussion in
\S III.C (super-first-order for smooth solutions, degrading toward first
order near the singularity).

\subsection*{C2b --- Non-convex star smoothing}
$r(\theta) = 0.25 + 0.10\cos(7\theta)$, $\varepsilon=0.05$, $T=0.3$.
Initial perimeter 3.332; final perimeter 1.217; \textbf{0.00\% of time
steps show a perimeter increase} (0 violations over 9\,000 steps),
matching Grayson's monotone-length-decrease.

\subsection*{C3 --- Topological merger}
Two disks of $R_0=0.15$ at $(\pm 0.3,0)$, $F=1$, $T=0.4$: components
$2 \rightarrow 1$; numerical merge time $0.1504$ vs analytical
$0.1500$; \textbf{relative error 0.27\%}. No detection or re-meshing.

\section{LLM-judge assessment}
The LLM judge (\texttt{argo:gpt-4o} at localhost:44497, free endpoint,
temperature 0) scored every sub-claim \texttt{pass} and returned
\texttt{overall\_verdict: REPLICATED} with the one-line summary
``All core claims of Osher \& Sethian (1988) are faithfully reproduced by
the independent numerical replication.''

\section{GENUINE CRITIQUE}
This is a deliberately scoped 2-D minimal replication; the following
limitations are real and should temper how far the result is extrapolated.

\begin{enumerate}
\item \textbf{Only 4 of 6 identified claims tested.} C4 (higher-order
ENO+RK order-of-accuracy test) and C5 (3-D / N-D generalization) are
explicitly out of scope. Claims about the method's high-order
extensions and its N-dimensional generality therefore rest on the
paper, not on our numbers.

\item \textbf{Convergence is short of the ideal.} The observed C2 order
is 1.3--1.9, not a clean 2. The paper explains this (singularity
formation as $R\to 0$), but a strict convergence check restricted to
early times / large $R$ would tighten the claim. We report the
composite order as-measured rather than a cherry-picked window.

\item \textbf{Merge-time error is a single number, not a convergence
sequence.} The 0.27\% figure at $N=251$ is agreement with the analytic
merge time, but we did not sweep the grid to confirm that this error
shrinks under refinement --- so the C3 quantitative result is
consistency at one resolution rather than a proven convergence.

\item \textbf{No comparison to an established level-set library.} The
implementation is checked against analytic solutions of highly
symmetric configurations (single circle, two symmetric disks) but not
against, e.g., \texttt{scikit-fmm}, \texttt{LSMLIB}, or FMM-based
reference codes. Bugs that happen to preserve the symmetric answers
would not be caught.

\item \textbf{No reinitialization / signed-distance repair.} The paper's
core scheme evolves $\phi$ directly, which is what we did; but any
production use would need periodic reinitialization to fight
$|\nabla\phi|$-drift, and we do not test that pathway.

\item \textbf{Boundary conditions differ from typical practice.}
Periodic BCs via \texttt{np.roll} keep the code trivial and work here
because the fronts stay well inside the domain, but they are not the
outflow / one-sided BCs commonly used in level-set applications and
would need to change for interfaces that touch the domain boundary.

\item \textbf{Explicit forward Euler only.} No RK, no implicit
treatment of the parabolic curvature term. The parabolic CFL
$\Delta t \lesssim \Delta x^{2}/(C\varepsilon)$ forces very small steps
at high resolution (up to $12{,}500$ steps at $N=301$), which is
acceptable at 2-D laptop scale but would be a real cost in 3-D
production. IMEX or implicit-curvature schemes would remove this cost
but were not implemented or compared.

\item \textbf{LLM judge is not an independent check.} The judge sees
only the paper's claims and our own numerical results; it cannot
detect an implementation bug that returns internally-consistent but
physically-wrong numbers. Its \texttt{pass} verdicts should be read as
``the numbers we produced are consistent with the paper's claims,''
not as independent validation.
\end{enumerate}

None of these limitations contradicts the REPLICATED verdict --- the
paper's four minimal, testable claims match analytic solutions to
$\lesssim 0.5\%$ --- but they do bound its scope.

\section{Verdict}
\textbf{REPLICATED.} The core mathematical claim of the paper --- that
recasting front propagation as a Hamilton--Jacobi PDE on an implicit
function and solving it with upwind Godunov + central-curvature schemes
gives a robust method that captures constant-speed motion,
curvature-driven motion, and topological changes on a fixed Eulerian
grid --- is independently reproduced from a from-scratch NumPy
implementation guided only by the equations in the PDF, with no
external level-set library.

\end{document}
